\documentclass[11pt]{article}
\usepackage[margin=1in]{geometry}
\usepackage{graphicx}
\usepackage{booktabs}
\usepackage{amssymb}
\usepackage{caption}
\captionsetup{font=small,labelfont=bf}

\newcommand{\figpic}[2]{%
  \begin{center}
  \includegraphics[width=0.62\textwidth]{#1}
  \vspace{8pt}
  {\small #2}
  \end{center}
  \vspace{12pt}}

\title{GSC Summary: Recovered Parameters and Model Results}
\author{}
\date{\today}

\begin{document}
\maketitle

\noindent Small figures so the whole project can be talked through. Under each
figure is a plain-language read of whether the result looks right and, if not,
why.

\section*{1. Recovered parameters (1965--2014)}

\figpic{../analysis/results/figures/consumption_demand_shifters.pdf}
  {This looks right. Services demand shifts up in both countries while demand
  for manufacturing and agriculture is flat to down, which is the standard
  structural-change pattern the model needs to reproduce.}

\figpic{../analysis/results/figures/trade_costs.pdf}
  {This looks right for the Rest of the World and flat for the US. Opening trade
  costs fall substantially, so the globalization channel is present. But the
  US-facing costs barely move, which is already a warning that freezing trade
  costs will not do much for the US.}

\figpic{../analysis/results/figures/productivities.pdf}
  {This looks right, with a caveat. Both countries grow in every sector and the
  Rest of the World catches up, which is the story we want. But the ROW rise is
  very large and worth checking against the raw capital and output series
  before relying on it.}

\figpic{../analysis/results/figures/financial_friction.pdf}
  {This looks right. The friction is small and gentle, staying close to one in
  both countries, which is what a wedge that only has to rationalize the
  observed net foreign asset build-up should look like.}

\figpic{../analysis/results/figures/investment_efficiency.pdf}
  {This needs a caveat. The tax is now recovered from the capital Euler
  equation rather than the law of motion, and it rises by roughly a factor of
  four over the sample. The drift is the model straining: nominal spending per
  worker grows faster than the discount factor and the measured rental rate can
  deliver, and the tax absorbs the difference.}

\figpic{../analysis/results/figures/investment_demand_shifters.pdf}
  {This looks right. The shifters move modestly and are dominated by services,
  and because they are read off observed investment the gap no longer falls on
  the consumption shifters. Nothing looks off.}

\section*{2. Model output and counterfactuals}

\noindent The counterfactual freezes bilateral trade costs at their 1965 values
and lets everything else move as recovered. Two runs: with the financial
friction active (headline) and with it switched off (robustness). The table
shows how much of the manufacturing value added decline each run explains over
the full sample and over the 1970--2007 window the paper reports.

\begin{table}[htbp]
\centering
\caption{Trade-cost counterfactual: change in value added shares (pp)}
\begin{tabular}{lcccc}
\toprule
 & \multicolumn{2}{c}{1965--2014} & \multicolumn{2}{c}{1970--2007} \\
\cmidrule(lr){2-3} \cmidrule(lr){4-5}
 & Man. & Goods & Man. & Goods \\
\midrule
Data                             & $-13.2$ & $-14.8$ & $ -9.9$ & $-11.4$ \\
Trade costs frozen at 1965       & $-12.0$ & $-13.5$ & $ -7.7$ & $ -9.1$ \\
\quad share of decline explained & $9.0\%$ & $8.5\%$ & $22.5\%$ & $19.8\%$ \\
Trade costs frozen, no friction  & $-14.7$ & $-16.7$ & $ -9.4$ & $-11.2$ \\
\quad share of decline explained & $-11.6\%$ & $-12.9\%$ & $5.4\%$ & $2.1\%$ \\
\bottomrule
\end{tabular}
\end{table}

\figpic{../analysis/results/figures/value_added_shares.pdf}
  {Look right mechanically. The model's runs hug the data, so it reproduces
  the structural change it is being asked to explain.}

\figpic{../analysis/results/figures/value_added_gap.pdf}
  {This needs more thought. Freezing trade costs removes only a small part of
  the manufacturing decline, so the headline channel explains little -- and
  how much it explains depends heavily on which years we look at.}

\figpic{../analysis/results/figures/trade_imbalances.pdf}
  {This looks right. With the friction the model traces the observed imbalance
  path; without it the path is basically flat, so this is the channel doing the
  work.}

\figpic{../analysis/results/figures/tradecost.pdf}
  {This looks right and confirms the trade-cost figure above: costs fall for
  the Rest of the World and stay put for the US, so freezing them changes
  little for the American economy.}

\medskip

\noindent \textbf{Net reading.} The model works, but the headline result is
weak. Trade costs account for at most a quarter of the manufacturing decline,
and the direction can even flip depending on the window and whether the
friction is included.

\section*{3. Outstanding questions}

\begin{enumerate}
\item \textbf{Do tradeable services count as goods?} The baseline keeps
wholesale, retail and hotels/restaurants in services and includes government
consumption. Moving wholesale and retail (and hotels and restaurants, which
the pre-2000 tables cannot separate) into goods and dropping government shifts
a large share of US consumption out of services and into manufacturing,
bringing services closer to the share the national accounts report. Under that
definition the recovered demand shifters change ordering -- manufacturing
starts above services and the two converge rather than diverge, matching the
2017 draft. Decide which sector definition is the right one.

\item \textbf{Things that may still be broken, in no particular order.}
\begin{itemize}
  \item The trade-cost answer is weak and window-dependent, and can flip sign
  when the friction is switched off. That may be a real finding, but it should
  be understood before it goes in the paper -- is it the flat US-facing costs,
  or the friction interacting with the frozen costs?
  \item The Rest of the World productivity rise is very large over fifty
  years; check it is not an artifact of the capital splice, the deflation, or
  the exchange-rate anchor.
  \item The recovered world interest rate is nearly flat. Confirm it matches
  the target we actually want to hit.
  \item After the exchange-rate fix the price-level ratio between the two
  regions moves in a reasonable band; recheck once more that the anchor is
  doing what it should for the level of trade costs.
\end{itemize}
\end{enumerate}

\end{document}